\section{Week 6: direct API trace boundary}

\subsection{Actual sequential trace}

\path{RawApiDirectObservedMu.ec} preserves the real execution order:
\begin{verbatim}
actual Sign raw ABI -> publish signature -> actual Verify raw ABI
       |                                      |
       v                                      v
SignRawApiMuTrace.run -> VerifyCryptolabMuTrace.run
\end{verbatim}
The same \texttt{sigu}, pre, and message arguments are used by both calls.
The exact theorem is
\begin{quote}
\small\path{raw_sign_then_verify_actual_exact_trace}.
\end{quote}
It proves equal result pairs and equal final memory. The traces add observations but retain the
entire signature core, verification residual computation, and rejection
behavior.

On the trace sequence,
\theoryname{sign_then_verify_trace_accept_implies_tail_reached} proves
\[
  \mathsf{VerifyAccept}\Longrightarrow\mathsf{TailRan},
\]
The companion binding lemma
\begin{center}
\small\path{sign_then_verify_accept_binds_observed_inputs}
\end{center}
derives both sets of VK/pre/message inputs from the post-state. No
\texttt{tail\_reached}
or observed hash value occurs in either precondition.

The intended full implication is
\[
 \mathsf{VerifyAccept}
 \Longrightarrow \mathsf{TailRan}
 \Longrightarrow
 \take_{32}(\mathsf{SignTrace.observed\_mu})
 = \mathsf{VerifyTrace.observed\_mu}.
\]

\subsection{What remains open}

The last implication is not yet compiled. The region-local theorem is a pRHL
relation between two hash procedure executions, while the desired result is a
unary postcondition about values stored by two calls that have already
completed inside one sequential trace. A product program, deterministic
functional hash specification, or justified replay rule is needed to
transport the former into the latter.

Repeating the Week 5 construction by placing both observations in a hash-call
precondition was rejected: it does not establish reachability from trace
execution. Adding an independent replay without proving equality to the
original observations was also rejected.

The remaining semantic step is named
\begin{center}
\small\path{OBL-MU-PRODUCT-REPLAY}.
\end{center}
The usual \texttt{byequiv} pattern does not directly convert the compiled
pRHL judgment into the required unary Hoare postcondition. Aligning the two
full traces relationally would leave the post-hash Sign rejection continuation
unpaired; eliminating it one-sidedly would require the still-open signing-loop
losslessness proof. That termination premise is outside this partial-
correctness seam and was not hidden.

Therefore:
\begin{itemize}
  \item \claimid{OBL-DIRECT-OBSERVED-MU}: \status{PARTIAL};
  \item \claimid{OBL-API-KEY-MEMORY-RAW-ACCEPT}: \status{PARTIAL};
  \item \(\delta_{\mu,\mathrm{raw\_api\_accept}}=0\): not claimed.
\end{itemize}

There is also no accepted Sign-output witness. That separate fact belongs to
the functional obligation
\begin{center}
\small\path{OBL-SIGN-OUTPUT-TAIL-REACH},
\end{center}
not to the accepted-forgery adapter.
